\chapter{Objetivos y metodología de trabajo}
\label{cap:objetivos}

Los objetivos delimitan tanto la construcción de la herramienta como la evaluación experimental de sus métodos de segmentación. Se separa la meta general de las tareas verificables que permiten alcanzarla.

\section{Objetivo general}

Diseñar, implementar y evaluar una herramienta de software basada en inteligencia artificial que, a partir de una fotografía de una hoja adhesiva impresa, rectifique la escena, segmente los \textit{toppers} mediante SAM 2 y genere trayectorias DXF registradas con el sistema de coordenadas del láser CNC, comparando su desempeño con Random Forest y una línea base clásica.

\section{Objetivos específicos}

\begin{description}
    \item[\textbf{OE1.}] Definir un procedimiento de captura con una cámara de consumo sobre una región delimitada de la mesa, adecuada para hojas adhesivas de formato A4.

    \item[\textbf{OE2.}] Implementar la detección de marcadores, la rectificación por homografía y la localización de una marca grabada por el láser para relacionar la imagen con el sistema de coordenadas de trabajo.

    \item[\textbf{OE3.}] Construir un conjunto sintético controlado y un protocolo con particiones independientes de entrenamiento, validación y prueba para evaluar métodos de segmentación.

    \item[\textbf{OE4.}] Entrenar y optimizar un Random Forest por píxel mediante validación agrupada, selección de hiperparámetros y ajuste del umbral sobre láminas completas.

    \item[\textbf{OE5.}] Seleccionar una estrategia de prompts para SAM 2 Hiera Tiny, evaluar su comportamiento en los dominios sintético y real e integrarlo como motor de segmentación de la aplicación.

    \item[\textbf{OE6.}] Implementar la extracción, simplificación y conversión de los contornos a coordenadas físicas, con exportación DXF condicionada a la existencia de un WCS válido.

    \item[\textbf{OE7.}] Comparar la línea base clásica, Random Forest y SAM 2 mediante Dice, IoU, coincidencia de bordes, error de componentes y tiempo, declarando con claridad las limitaciones de los datos reales y de la validación física.
\end{description}

\section{Metodología de trabajo}

El proyecto se organizó en siete fases, con un punto de control común: separar las decisiones de desarrollo de la evaluación final. La delimitación del problema y la revisión técnica dieron lugar a los requisitos y al montaje de captura. Una vez implementado el registro geométrico, desde los marcadores de la hoja hasta el sistema de coordenadas de trabajo, se construyó el conjunto sintético y se reunieron las capturas reales. Las particiones sintéticas quedaron fijadas antes del entrenamiento y de cualquier selección de parámetros, mientras que las trece capturas del taller se reservaron para la evaluación externa.

La comparación reunió una línea base de visión clásica, un Random Forest supervisado por píxel y SAM 2 como modelo fundacional condicionado mediante cajas. En Random Forest, el entrenamiento alimentó la búsqueda agrupada y las hojas de validación resolvieron la elección del candidato y el umbral. Sobre ese mismo subconjunto de validación se seleccionó el margen de caja de SAM 2. La prueba no intervino en esas decisiones ni en los parámetros clásicos fijados durante el desarrollo previo. La máscara de SAM 2 se conectó después con la extracción de contornos dentro de la herramienta, cuya exportación DXF queda bloqueada si falta una referencia WCS válida. La evaluación, el análisis de amenazas y la comprobación de reproducibilidad completaron el estudio.

El Capítulo~\ref{cap:metodologia} desarrolla con detalle el protocolo técnico, las unidades experimentales, la optimización de hiperparámetros, las métricas y las condiciones de prueba.
